\section{Preliminaries}\label{sec:prel}

% \commentt{Might remove some of the following if it appear in the intro} 
Unless stated explicitly otherwise, the graphs we consider are simple, so that in particular they do not contain any parallel edges. We denote the number of vertices in the graph by $n$ and the number of edges by $m$. Every vertex $v$ in the graph has a unique id, denoted $id(v)$, and its degree is denoted by $d(v)$.

We work in what is known as the \emph{general graph model}~\cite{ParnasR02,KKR04}. In particular, under this model,
the \emph{distance} of a graph $G$ to $C_k$-freeness, denoted $dist(G,C_k\text{-free})$, is the minimum fraction of edges that should be  removed from $G$ in order to obtain a $C_k$-free graph.
As for the allowed queries, 
a neighbor query to the $i$th neighbor of a vertex $v$ is denoted by $nbr(v,i)$, and to its degree by $deg(v)$. A pair query between two vertices $v_1$ and $v_2$ is denoted by $pair(v_1,v_2)$.
Given query access to a graph $G$ and a parameter $\eps$, a one-sided error testing algorithm for $C_k$-freeness should accept $G$ if it is $C_k$-free, and should reject $G$ with probability at least $2/3$ if $dist(G,C_k\text{-free}) > \eps$.
If the algorithm may also reject $C_k$-free graphs with probability at most $1/3$, then it has two-sided error.

As noted in the introduction, we assume our algorithms for graphs whose arboricity is not promised to be constant,  are given an upper bound $\alpha$ on the arboricity $arb(G)$ of the tested  graph $G$, and their complexity depends on this upper bound. Alternatively, if the algorithm is provided with the number of edges, $m$, then it may run a procedure from~\cite{Reut} to obtain a value $\alpha^*$ that with high constant probability satisfies the following: (1) $\alpha^* \leq 2arb(G)$; (2) The number of edges between vertices whose degree is at least $\alpha^*/(c\eps)$ for a constant $c$ is at most $(1-\eps/c')m$ (for another, sufficiently large, constant $c'$).
Up to polylogarithmic factors in $n$, the query complexity and running time of the procedure are $O(arb(G)/\eps^3)$ with high probability (assuming the average degree is $\Omega(1)$). 

%\commentt{Elaborate more on $\alpha^*$.}

Throughout this work we assume, whenever needed, that $\epsilon$ is upper bounded by some sufficiently small constant (or else it can be set to that constant).

